\documentclass[12pt]{article}
%% BUGS:
%% - The \begin{problem} ... \end{problem} is only used in exams??

\usepackage{url, graphicx}
\PassOptionsToPackage{normalem}{ulem} % magic
\usepackage{href-ul}

% page layout
\usepackage[letterpaper, textheight=9.5in, textwidth=5.5in]{geometry}
\usepackage{setspace}
\setstretch{1.08}
\pagestyle{plain}
\sloppy\sloppypar\raggedbottom\frenchspacing\thispagestyle{empty}

% problem set header
\newcommand{\psname}{Problem Set}
\newcommand{\startps}[1]{\section*{Computational Physics --- \psname~{#1}}}

% problem formatting
\usepackage{amsthm}
\theoremstyle{definition}
\newcommand{\problemname}{Problem}
\newtheorem{problem}{\problemname}
\newcommand{\probref}[1]{\problemname~\ref{#1}}
\theoremstyle{remark}
\newcommand{\notename}{Note}
\newtheorem{note}{\notename}[problem]
\newcommand{\noteref}[1]{\textit{\notename}~\ref{#1}}
\newcounter{subpart}[problem]
\renewcommand{\thesubpart}{\alph{subpart}}
\newcommand{\subpart}{\par\addvspace{0.75ex}\refstepcounter{subpart}\textsl{(\thesubpart)}~}
\newcommand{\subpartref}[1]{\textsl{(\ref{#1})}}

% words
\newcommand{\foreign}[1]{\textsl{#1}}
\newcommand{\vs}{\foreign{vs}}

% code
\newcommand{\code}[1]{\texttt{\detokenize{#1}}}
\usepackage{xcolor}
% Define custom colors for code highlighting
\definecolor{codegreen}{rgb}{0,0.6,0}
\definecolor{codegray}{rgb}{0.5,0.5,0.5}
\definecolor{codepurple}{rgb}{0.58,0,0.82}
\definecolor{backcolour}{rgb}{0.98,0.98,0.99}
\usepackage{listings}
\lstdefinestyle{pythonstyle}{
    backgroundcolor=\color{backcolour},   
    commentstyle=\color{codegreen},
    keywordstyle=\color{magenta},
    numberstyle=\tiny\color{codegray},
    stringstyle=\color{codepurple},
    tabsize=4
}
\lstset{
    basicstyle=\small\ttfamily,
    columns=fixed,
    breaklines=false,
    keepspaces=true,
    showspaces=false,
    showstringspaces=false,
    showtabs=false,
    showlines=*,
    style=pythonstyle,
    language=Python,
    literate={~}{{\url{~}}}1
}

% math
\usepackage{amsmath}
\newcommand{\dd}{\mathrm{d}}
\newcommand{\e}{\mathrm{e}}

% primary units
\newcommand{\rad}{\mathrm{rad}}
\newcommand{\kg}{\mathrm{kg}}
\newcommand{\m}{\mathrm{m}}
\newcommand{\s}{\mathrm{s}}

% secondary units
\renewcommand{\deg}{\mathrm{deg}}
\newcommand{\km}{\mathrm{km}}
\newcommand{\cm}{\mathrm{cm}}
\newcommand{\mm}{\mathrm{mm}}
\newcommand{\ft}{\mathrm{ft}}
\newcommand{\mi}{\mathrm{mi}}
\newcommand{\AU}{\mathrm{a.u.}}
\newcommand{\au}{\mathrm{a.u.}}
\newcommand{\ns}{\mathrm{ns}}
\newcommand{\h}{\mathrm{h}}
\newcommand{\yr}{\mathrm{yr}}
\newcommand{\N}{\mathrm{N}}
\newcommand{\J}{\mathrm{J}}
\newcommand{\eV}{\mathrm{eV}}
\newcommand{\W}{\mathrm{W}}
\newcommand{\Pa}{\mathrm{Pa}}

% derived units
\newcommand{\mps}{\m\,\s^{-1}}
\newcommand{\mph}{\mi\,\h^{-1}}
\newcommand{\mpss}{\m\,\s^{-2}}

% draft comments
\newcommand{\hoggsay}[1]{{\color{blue}\textbf{#1}}}

\begin{document}

\startps{8}
Due Friday 2026 October 30 at 18:00 on Brightspace.
Be sure to explicitly give credit where it is due, on all problems.
Also credit buddies, web pages, and LLMs (as per course policy).
Failure to give explicit credits will result in points deducted.

\begin{problem}
  We're going to do a modification of Exercise~8.7 in Newman, \textit{Computational Physics, 2ed},
  for a golf ball (not a cannonball).
  Look up the properties of a golf ball and use these instead of what's written in that problem.
  Once you have done all parts \textsl{(a)} through \textsl{(c)}, do one final part
  \setcounter{subpart}{3}
  \subpart The fastest any mechanical device could ever launch a golf ball is about one half the speed of sound in air (``mach 0.5'').
  Any faster and the ball would ablate.
  Play around with your integrator to see how far a golfer could hit a golf ball if they could hit the ball at half the speed of sound in air.
  Comment on your answer in the context of the longest golf drives ever hit.
\end{problem}

\begin{note}
  Use your leapfrog integrator, or use any standard Python integrator.
\end{note}

\begin{note}
  On the last part, you could use ideas from optimization to find the
  farthest the ball can fly over level ground when launched at mach
  0.5.
\end{note}

\begin{problem}
  You can think of an ODE as a map:
  You start with initial conditions $p,q$ and you integrate forward in time to get a new state $p,q$.
  That's a mapping $p,q\rightarrow p,q$.
  The simplest (maybe) interesting map is the \emph{logistic map}:
  \begin{align}
    x &\leftarrow r\,x\,[1 - x] ~,
  \end{align}
  where $r$ is a continuous parameter setting how strong the nonlinearity is.

  Write simple code that iterates this map.
  At every $r$ value in \code{np.arange(2.5, 4.0, 0.001)} do the following:
  Starting at $x=0.5$, iterate the map 1000 times as a ``burn-in phase'' at one fixed value of $r$.
  Then make a scatter plot of $x$ \foreign{vs} $r$ for the next 100 iterations.
  Then move to the next $r$ value.

  Look up this map on the internets.
  What does it have to do with chaos?
\end{problem}

\begin{note}
  Play around with your plotting so that you can see all the amazing structure.
  You will want to make your points tiny!
\end{note}

\begin{note}
  Most values of $r$ prodice one fixed point, or 2, or $2^n$, or none.
  The chaotic transitions happen when triples of fixed points appear.
  Can you find an $r$ value that has a triplet of fixed points?
  That's called a 3-cycle, and it is critical to the onset of chaos.
\end{note}

\begin{note}
  Hogg speaks of ``moving to the next $r$ value'' but if you are clever,
  you can do all $r$ values in your list all at once, in parallel.
  If you want to think about how to code, figure out how to do this problem
  with only one for loop (the loop over iterations of the map, which requires
  sequential execution).
\end{note}

\begin{problem}
  \hoggsay{Surface of section.}
\end{problem}

\end{document}
